\documentclass{amsart}

\input{../preamble_general.tex}

\setlength{\textheight}{700pt}

\author{Jalen Chrysos}
\title{Embedded Graphs}

\begin{document}
	\maketitle 
	We've talked about planar graphs, graphs that can be embedded in $\R^2$ without edges crossing. But what if we replace $\R^2$ by some other surface $S$, for example a torus?\\
%	We've learned that if a graph $G$ is embedded in the plane, then it satisfies Euler's formula:
%	$$
%	V - E + F = 2
%	$$
%	where $V$ is the number of vertices of $G$, $E$ is the number of edges, and $F$ is the number of faces. We've also seen that $G$ can be embedded in the plane if and only if $G$ does not contain $K_5$ or $K_{3,3}$ as minors.\\
%	
%	\underline{\textbf{Question}}: For a given topological space $S$, what can we say about graphs embedded in $S$?\\
	
	\noindent Technical note: we only consider \textit{cellular embeddings}, which are embeddings of $G$ such that the faces of $G$ are homeomorphic to disks---informally, they have no holes.\\
%	
%	$$\textbf{DIAGRAM: CELLULAR VS NONCELLULAR EMBEDDING}$$
	
	\noindent 
	If $G$ has a cellular embedding into $S$, it may not be that $V-E+F=2$, but it is true that $V-E+F$ is some constant depending on $S$ (and not on $G$), which we denote $\chi(S)$:
	$$
	V - E + F = \chi(S).
	$$
	$\chi(S)$ is called the \textbf{\textit{Euler characteristic}} of $S$. In practice it can be computed just by knowing this fact and picking a fairly simple graph with a cellular embedding into $S$.\\
	
	\noindent \textbf{Problem 1}: Compute the Euler characteristics of the following surfaces:
	\begin{enumerate}[(i)]
		\item Plane
		\item Sphere
		\item Torus
		\item $2$-holed torus\\
	\end{enumerate}
	
	\noindent \textbf{Problem 2}: Suppose a simple graph $G$ has a cellular embedding into $S$.
	\begin{enumerate}[(a)]
		\item Show that $E\leq 3(V-\chi(S))$.
		\item Suppose that the shortest cycle in $G$ has length $k$. Show that 
		$$
		E \leq \frac{k}{k-2} (V-\chi(S)).
		$$
	\end{enumerate}
	\vspace{0.1cm}
	
	\noindent \textbf{Problem 3}: For which $n$ does $K_n$ have a cellular embedding into the torus?\\
	
	\begin{center}
	\begin{tabular}{|c|c|c|c|c|c|c|c|}
		\hline 
		$n$ & $\;3\;$ & $\;4\;$ & $\;5\;$ & $\;6\;$ & $\;7\;$ & $\;8\;$ & \dots \\
		\hline 
		Embedding? & & & & & & & \\[3pt]
		\hline 
	\end{tabular}
	\end{center}
	\vspace{0.3cm}
	
	\noindent \textbf{Problem 4}: For which $m,n$ does $K_{m,n}$ have a cellular embedding into the torus?\\

	\begin{center}
		\begin{tabular}{|c|c|c|c|c|c|c|c|}
			\hline 
			$m,n$ & $\; 2\; $ & $\;3\;$ & $\;4\;$ & $\;5\;$ & $\;6\;$ & $\;7\;$ & \dots \\[3pt]
			\hline 
			$2$ & & & & & & & \\[3pt]
			\hline 
			$3$ & & & & & & & \\[3pt]
			\hline 
			$4$ & & & & & & & \\[3pt]
			\hline 
			$5$ & & & & & & & \\[3pt]
			\hline 
			$6$ & & & & & & & \\[3pt]
			\hline 
			$7$ & & & & & & & \\[3pt]
			\hline 
			$\vdots$ & & & & & & & \\
			\hline 
		\end{tabular}
	\end{center}
	
	\vspace{0.3cm}
	
	Robertson and Seymour proved in 2004 that for any surface $S$, there is some finite set $\G(S)$ of graphs such that $G$ can be cellularly embedded into $S$ if and only if it has no minor in $\G(S)$. For planar graphs, as we saw, $\G(\R^2)=\{K_5,K_{3,3}\}$. For the torus, $\G(T)$ is still not fully known, but it is known that it contains at least 17,000 graphs.
	
\end{document}